\subsection{Comparison with the classical \texorpdfstring{$r$}{r}-matrix construction}
\label{sec:class5-rmatrix-comparison}

The time Lax matrix in Eq.~\eqref{eq:class5-time-Lax-final} was obtained by
starting from the finite-lattice time operator.  We now construct the time
component independently from the classical $r$-matrix and monodromy of the
continuum model.  This provides a check that does not use the finite-lattice
time operator $A_{a,n}$.

To keep the two roles of the spectral parameter separate, let
$\zeta\in\mathbb C$ denote the \emph{generating spectral parameter} used to
expand the transfer matrix, while $\lambda$ continues to denote the Lax
spectral parameter of the final pair $(U_5,V_5)$.  Define the rank-one
coherent-state projector
\begin{equation}
\rho(x,t)
:=
\frac12
\begin{pmatrix}
1+S^3&S^-\\
S^+&1-S^3
\end{pmatrix},
\qquad
\rho^2=\rho,
\qquad
\tr\rho=1,
\label{eq:class5-rho-definition}
\end{equation}
and the deformation matrix in the auxiliary space
\begin{equation}
K^{\mathrm{cl}}_5(x,t)
:=
\frac12
\begin{pmatrix}
S^+&-(1+S^3)\\
0&0
\end{pmatrix}.
\label{eq:class5-Kaux-definition}
\end{equation}
These are precisely the lower symbols in
Eqs.~\eqref{eq:class5-P-lower-symbol} and
\eqref{eq:class5-K-lower-symbol}.  The spatial Lax matrix therefore admits the
form
\begin{equation}
U_5(x,t;\lambda)
=
\frac{\rho(x,t)}{2\lambda}
+\alpha_5 K^{\mathrm{cl}}_5(x,t).
\label{eq:class5-U-rho-K}
\end{equation}

We first fix the classical $r$-matrix in the conventions of this paper.  Let
$P_{12}$ be the permutation operator on two copies of the two-dimensional
auxiliary space.  We also define
\begin{equation}
P_{\uparrow}:=
\begin{pmatrix}1&0\\0&0\end{pmatrix},
\qquad
E_{+}:=
\begin{pmatrix}0&1\\0&0\end{pmatrix},
\label{eq:class5-rmatrix-elementary-matrices}
\end{equation}
and the constant Jordanian tensor
\begin{equation}
C_{12}
:=
P_{\uparrow}\otimes E_{+}
-
E_{+}\otimes P_{\uparrow}.
\label{eq:class5-jordanian-tensor}
\end{equation}
A direct calculation with the spin Poisson bracket
\eqref{eq:class5-poisson-bracket} gives
\begin{equation}
r^{(5)}_{12}(\zeta,\lambda)
=
\frac{\ii}{2(\zeta-\lambda)}P_{12}
+\ii\alpha_5 C_{12},
\label{eq:class5-classical-rmatrix}
\end{equation}
for which
\begin{align}
\{U_{5,1}(x;\zeta),U_{5,2}(y;\lambda)\}_5
={}&
\bigl[
 r^{(5)}_{12}(\zeta,\lambda),
 U_{5,1}(x;\zeta)+U_{5,2}(y;\lambda)
\bigr]\delta(x-y).
\label{eq:class5-linear-rmatrix-algebra}
\end{align}
Here
$U_{5,1}:=U_5\otimes\Id_2$ and
$U_{5,2}:=\Id_2\otimes U_5$ act on the first and second auxiliary spaces,
respectively.  The rational part of
Eq.~\eqref{eq:class5-classical-rmatrix} is the usual isotropic
Landau--Lifshitz contribution, whereas the constant term contains the Class 5
deformation.  This structure is consistent with the classical $r$-matrix
obtained from the Class 5 quantum model in
Ref.~\cite{deLeeuwFontanellaNietoGarcia2026}.

For completeness, define the classical monodromy matrix by
\begin{equation}
T(x,y;\zeta)
:=
\mathcal P\exp\!\left(
\int_y^x U_5(s,t;\zeta)\,ds
\right),
\label{eq:class5-classical-monodromy}
\end{equation}
and the associated transfer matrix on the periodic interval of length $\ell$
by
\begin{equation}
\tau(\zeta):=\tr T(\ell,0;\zeta).
\label{eq:class5-classical-transfer}
\end{equation}
For a linear classical $r$-matrix algebra of the form
\eqref{eq:class5-linear-rmatrix-algebra}, the standard monodromy construction
gives the generating time matrix
\begin{equation}
\mathbb V_2(x;\zeta,\lambda)
=
\frac{1}{\tau(\zeta)}
\tr_1\!\left[
T_1(\ell,x;\zeta)
\,r^{(5)}_{12}(\zeta,\lambda)\,
T_1(x,0;\zeta)
\right],
\label{eq:class5-rmatrix-generating-V}
\end{equation}
where the trace is taken over the first auxiliary space
\cite{SemenovTianShansky1983,DoikouKaraiskos2011}.

The pole of Eq.~\eqref{eq:class5-U-rho-K} at small $\zeta$ is proportional to
the rank-one projector $\rho$.  The local expansion of the monodromy is
therefore governed by a rank-one projector $\Pi(x;\zeta)$ satisfying
\begin{equation}
\Pi^2=\Pi,
\qquad
\tr\Pi=1,
\qquad
\partial_x\Pi=[U_5(x;\zeta),\Pi],
\label{eq:class5-local-projector-equations}
\end{equation}
with the branch
\begin{equation}
\Pi(x;\zeta)
=
\rho(x)+\zeta\Pi_1(x)+\mathcal O(\zeta^2).
\label{eq:class5-local-projector-expansion}
\end{equation}
Expanding Eq.~\eqref{eq:class5-local-projector-equations} to first order fixes
the coefficient uniquely as
\begin{equation}
\Pi_1
=
2[\rho,\partial_x\rho]
-2\alpha_5[\rho,[K^{\mathrm{cl}}_5,\rho]].
\label{eq:class5-Pi1-solution}
\end{equation}
Thus the first correction contains both the ordinary spatial-gradient term
and the deformation term.

Using Eq.~\eqref{eq:class5-local-projector-expansion} in the generating
matrix \eqref{eq:class5-rmatrix-generating-V}, write
\begin{equation}
\mathbb V_2(x;\zeta,\lambda)
=
\mathbb V^{(0)}(x;\lambda)
+\zeta\mathbb V^{(1)}(x;\lambda)
+\mathcal O(\zeta^2).
\label{eq:class5-generating-V-expansion}
\end{equation}
The coefficient relevant for the Class 5 Hamiltonian flow is
\begin{align}
\mathbb V^{(1)}
={}&
-\frac{\ii}{2\lambda}
\left(
\Pi_1+\frac{\rho}{\lambda}
\right)
\nonumber\\
&+\ii\alpha_5\left[
\tr(\Pi_1P_{\uparrow})E_{+}
-\tr(\Pi_1E_{+})P_{\uparrow}
\right].
\label{eq:class5-generating-V-first-coefficient}
\end{align}
The standard monodromy formula is conventionally written with the Hamiltonian
in the first argument of the Poisson bracket.  Our Hamiltonian evolution in
Eq.~\eqref{eq:class5-hamiltonian-flow-definition} uses the opposite order.
With this convention, the time Lax matrix obtained from the classical
$r$-matrix route is therefore
\begin{equation}
V_{5,r}(x;\lambda):=-\mathbb V^{(1)}(x;\lambda).
\label{eq:class5-rmatrix-time-V-definition}
\end{equation}

Substituting Eq.~\eqref{eq:class5-Pi1-solution} into
Eq.~\eqref{eq:class5-generating-V-first-coefficient} and using only the unit
spin constraint gives
\begin{equation}
V_{5,r}(x;\lambda)-V_5(x;\lambda)
=
\frac{\ii}{4\lambda^2}\Id_2.
\label{eq:class5-rmatrix-vs-quantum-V}
\end{equation}
The matrix on the right-hand side is independent of the spin field and of the
spatial coordinate.  Its spatial derivative vanishes and it commutes with
$U_5$.  Hence $V_{5,r}$ and the finite-lattice result $V_5$ generate the same
zero-curvature equation.  The two constructions thus recover the same Class
5 time Lax component up to the scalar freedom already encountered in
Eq.~\eqref{eq:class5-DeltaV}.

